\chapter{激进的游戏}

\section{本章导入}

德州扑克不是被动等待好牌的游戏。很多新手刚开始打牌时，最自然的策略是等待强牌：拿到好牌就下注，没牌就过牌或弃牌；成牌就继续，没中就放弃。这样的打法表面上稳健，实际上会让牌手失去一半赢底池的方式。

获得底池通常有两种方法：

\begin{enumerate}
\item 摊牌时拥有最好牌；
\item 通过进攻让对手弃牌。
\end{enumerate}

新手往往只会第一种方式。他们必须等自己成牌，必须等自己领先，必须等公共牌帮助自己，才敢投入筹码。这样一来，很多原本可以通过主动进攻赢下的底池，都被他们放弃了。更重要的是，对手会很容易读懂他们：下注就是有牌，过牌就是放弃。这样的玩家很难给对手制造持续压力。

优秀牌手不只是等好牌，而是在合适的局面主动争取底池。他们理解主动权，理解弃牌权益，理解什么牌适合半诈唬，理解什么牌面适合高频进攻，也理解什么对手可以被施压。真正的激进不是乱 bluff，而是有条件、有目标、有计划地施压。

激进并不等于鲁莽。鲁莽是没有范围优势也进攻，没有牌面支持也进攻，对手不会弃牌也进攻，自己的牌没有功能也进攻。有效激进则不同，它建立在三个条件之上：

\begin{center}
\textbf{我的范围有优势。\\
我的牌有合适的功能。\\
对手有足够多可以弃掉的牌。}
\end{center}

这就是本章的核心框架。

本章要解决的问题包括：为什么不能只等好牌？什么是弃牌权益？激进和鲁莽有什么区别？什么情况下应该主动抢夺底池？什么情况下激进反而是亏损？加注除了代表强牌之外，还有哪些意义？

学完本章后，你应该逐渐从“只有牌才敢打”的被动玩家，转向“在合适条件下主动争取底池”的思考型玩家。

\section{核心概念}

\subsection{为什么扑克需要激进性}

德州扑克中，如果你只靠摊牌赢底池，就会天然处于劣势。原因很简单：多数时候，你并不会在翻牌、转牌或河牌都拿到明显强牌。很多手牌会错过牌面，很多成牌只是中等牌力，很多听牌最终不会完成。如果你只能在自己牌力很强时才下注，那么你的赢池方式会非常有限。

激进性的重要意义在于，它让你拥有第二种赢法：让对手弃牌。

例如，你在按钮位 open，大盲跟注。翻牌为 \hand{A♠ 7♣ 2♦ rainbow}。你手持 \hand{Q♠ J♣}，没有成牌。如果你只看自己的牌，你会选择过牌放弃。但这个牌面对你的按钮 open range 更有利，大盲有许多没有击中的牌。你的小尺度持续下注可能直接赢下底池。此时你并不是靠摊牌赢，而是靠范围优势和弃牌权益赢。

如果一个玩家完全缺乏激进性，对手可以轻松应对他：

\begin{itemize}
\item 他下注时，对手知道他通常有牌；
\item 他过牌时，对手可以频繁下注抢夺；
\item 他拿中等牌力时，经常被迫猜测；
\item 他错过牌面时，几乎自动放弃；
\item 他很难从对手的弃牌中获利。
\end{itemize}

扑克中的盈利来自两个方向：一是强牌拿价值，二是让对手在不舒服的范围中弃牌。没有激进性，就相当于只用一条腿走路。

但激进性必须建立在理性基础上。有效激进不是更多下注，而是更准确地下注；不是更多 bluff，而是选择更好的 bluff；不是永远进攻，而是在条件有利时进攻。

\subsection{主动权}

主动权通常来自翻前或前一街的加注、下注行动。拥有主动权的一方，在后续街更容易代表强范围，也更容易通过持续下注给对手施压。

例如你翻前 open raise，大盲跟注。你是翻前进攻方，翻后许多牌面上你可以继续下注。对手作为防守方，必须面对你的下注决定是否继续。这种结构本身就给了你主动权。

主动权的价值主要体现在三个方面。

第一，它让你拥有先发压力。你可以通过下注迫使对手筛选范围，而不是被动等待对手下注。

第二，它让你的范围故事更连贯。翻前加注、翻牌下注、转牌继续，这些行动可以共同表达强范围。如果牌面支持你的范围，对手的中等牌力会承受越来越大压力。

第三，它让你拥有更多赢池方式。你可以摊牌赢，也可以通过下注让对手弃牌；而被动玩家更多只能依赖摊牌。

不过，主动权不是无限开火的许可证。你是翻前加注者，不代表每个翻牌都必须 C-bet；你翻牌下注了，不代表每个转牌都必须开第二枪。主动权只是给你进攻资格，是否进攻仍然要看范围优势、牌面结构、手牌功能和对手类型。

\subsection{弃牌权益}

弃牌权益，指的是你通过下注或加注让对手弃牌，从而直接赢下底池的可能性。它是激进策略的核心。

如果你下注后，对手有一定概率弃牌，那么即使你的牌当前不是最好牌，你仍然可以通过下注盈利。尤其当你的牌同时拥有一定牌面权益时，下注会变得更有价值。

例如你持有同花听牌。你下注后，对手可能弃牌，你直接赢底池；如果对手跟注，你仍然有机会在后续街完成同花。这就是弃牌权益和牌面权益叠加带来的力量。

弃牌权益取决于几个因素：

\begin{itemize}
\item 对手范围中有多少弱牌和中等牌力；
\item 你的下注尺度是否足够制造压力；
\item 你的行动线是否可信；
\item 牌面是否支持你代表强牌；
\item 对手是否有弃牌能力。
\end{itemize}

新手经常忽略弃牌权益，所以他们会觉得“我现在没有牌，下注就是乱打”。但事实上，很多下注并不是因为你当前已经领先，而是因为你可以让对手放弃更好牌，或者放弃仍有权益的牌。

不过，弃牌权益不是幻想出来的。对不会弃牌的玩家，弃牌权益很低。面对跟注站，诈唬和半诈唬价值下降，价值下注价值上升。激进之前，必须判断对手是否真的会弃牌。

\subsection{激进和鲁莽的区别}

激进和鲁莽表面上都表现为下注、加注、施压，但本质完全不同。

激进是有条件的进攻。你进攻之前，知道自己为什么下注，知道自己代表什么范围，知道对手有哪些牌会弃，知道如果被跟注或加注应该怎么办。激进的目标是让对手犯错：用更差牌付钱，或者弃掉本该继续的权益。

鲁莽是没有条件的进攻。牌面不支持自己也打，对手不弃牌也打，手牌没有权益也打，被加注后没有计划也打。鲁莽的目标往往不是长期 EV，而是情绪释放、证明自己、追回损失或不想显得软弱。

可以用以下表格区分两者：

\begin{center}
\begin{tabular}{lll}
\toprule
\textbf{比较维度} & \textbf{有效激进} & \textbf{鲁莽进攻} \\
\midrule
范围基础 & 有范围优势或坚果优势 & 不考虑范围 \\
手牌功能 & 价值牌、强听牌、好诈唬候选 & 任意牌硬打 \\
对手类型 & 对手有可弃牌范围 & 对跟注站强行 bluff \\
下注目的 & 价值或诈唬明确 & 试探、发泄、追损 \\
后续计划 & 知道转牌、河牌如何继续 & 被跟注后大脑空白 \\
\bottomrule
\end{tabular}
\end{center}

所以，本章强调激进，并不是鼓励新手乱打。恰恰相反，真正的激进比被动更需要纪律。因为你投入筹码更多，错误成本也更高。如果没有条件和计划，所谓激进只会变成更快输钱。

\subsection{进攻不是乱 bluff，而是有条件地施压}

很多新手一听到“激进”，就以为是多诈唬、多开火、多加注。其实进攻并不等于乱 bluff。进攻的本质是施压，而施压必须建立在对手承受压力的能力有限这一事实之上。

当对手范围中有大量中等牌力时，进攻最有效。因为中等牌力最怕面对多街压力。它们能击败一些弱牌和诈唬，但无法舒服地承受大底池。如果你的下注线能够代表强牌，对手就会被迫弃掉一部分仍然有摊牌价值的牌。

例如，对手在大盲跟注你的按钮 open。翻牌 \hand{A♠ 7♣ 2♦ rainbow}，他 check-call 小注。转牌 \hand{K}，他再次过牌。这个 K 强化你的范围。你继续下注，可以让他的 7x、小对子、弱 A 甚至部分 Kx 承受压力。这里的进攻不是乱 bluff，而是利用范围优势和转牌优势施压。

但如果对手范围非常强，或者牌面明显击中对手，强行进攻就会变差。例如你在前位 open，大盲跟注，翻牌 \hand{9♠ 8♠ 7♦ two-tone}，转牌 \hand{6}。这类牌面极大强化大盲范围。你如果拿空气牌继续大额下注，往往不是有效激进，而是鲁莽。

进攻前要问：

\begin{center}
\textbf{对手有哪些牌会因为我的下注而难受？}
\end{center}

如果对手大多是强牌，不会难受；如果对手是跟注站，不愿弃牌；如果你的故事不可信，那么进攻价值就很低。

\subsection{什么牌适合半诈唬}

半诈唬是激进游戏中最重要的工具之一。它介于价值下注和纯诈唬之间：你当前未必领先，但如果对手弃牌，你直接赢；如果对手跟注，你仍然有机会改善成强牌。

适合半诈唬的牌通常具备以下特点：

\begin{itemize}
\item 有较强牌面权益，例如坚果同花听牌、两头顺听牌、组合听牌；
\item 有额外高张价值，例如同花听牌加两个高张；
\item 有后门权益，例如后门同花加后门顺子；
\item 有阻断效果，例如持有关键 A 或 K，减少对手强牌组合；
\item 被跟注后仍然有后续计划。
\end{itemize}

例如，你持有 \hand{A♠ Q♠}，翻牌为 \hand{J♠ T♣ 4♦}，其中两张与你同花。这是非常好的半诈唬候选。你有坚果同花听牌、两张高牌以及顺子潜力。下注后，对手可能弃掉一些小对子或弱牌；即使被跟注，你也有许多好转牌。

相反，弱卡顺、低同花听牌、没有额外权益的低质量听牌，并不适合频繁大规模半诈唬。它们被跟注后权益不足，完成后也可能被更强牌压制。

半诈唬不是“有听牌就打”，而是“有足够牌面权益和弃牌权益时打”。

\subsection{什么牌面适合高频进攻}

并不是所有牌面都适合高频进攻。牌面是否适合进攻，取决于它对谁的范围更有利，以及它是否容易让对手范围错过。

适合高频进攻的常见牌面包括：

\begin{itemize}
\item 高牌干燥面，例如 \hand{A♠ 7♣ 2♦ rainbow}、\hand{K♠ 8♣ 3♦ rainbow}；
\item 明显有利于翻前加注者的牌面；
\item 对手防守范围中有大量没有击中的牌；
\item 听牌较少、对手难以强力反击的牌面；
\item 转牌或河牌强化你范围的牌面。
\end{itemize}

例如你翻前 open，大盲跟注。翻牌 \hand{A♠ 7♣ 2♦ rainbow}，这个牌面对你的 open range 通常更有利。你可以用较宽范围小注进攻，包括强 A、弱 A、高对、部分空气牌和后门权益牌。

不适合盲目高频进攻的牌面包括：

\begin{itemize}
\item 中低连接面，例如 \hand{9♠ 8♣ 7♦}、\hand{8♠ 7♣ 6♦}、\hand{T♠ 9♣ 8♦}；
\item 两张同花且顺子听牌丰富的湿润面；
\item 更容易击中大盲防守范围的牌面；
\item 对手拥有更多两对、暗三条和顺子组合的牌面；
\item 多人底池中的复杂牌面。
\end{itemize}

在这些牌面上，进攻方需要更谨慎。即使你是翻前加注者，也不能自动持续下注。

\subsection{什么对手适合被施压}

激进策略必须考虑对手类型。不是所有对手都适合被施压，也不是所有对手都应该被频繁诈唬。

适合被施压的对手通常有以下特点：

\begin{itemize}
\item 翻前入池较宽，翻后范围较弱；
\item 翻牌喜欢跟注，但转牌和河牌面对压力会弃牌；
\item 害怕大底池；
\item 中等牌力不愿意承受多街下注；
\item 能够理解牌面威胁，并愿意弃牌。
\end{itemize}

这类对手经常在翻牌跟得较宽，但转牌或河牌面对大注会弃掉大量中等牌力。对他们，合理的第二枪和第三枪会有较好效果。

不适合频繁施压的对手包括：

\begin{itemize}
\item 跟注站：有任何对子都想看摊牌；
\item 极端激进玩家：频繁反加，喜欢对抗；
\item 范围很紧的玩家：继续到后面街时通常很强；
\item 不理解牌面威胁的娱乐玩家：不会因为牌面变危险而弃牌。
\end{itemize}

面对跟注站，最好的策略不是复杂诈唬，而是用价值牌多下注。面对过紧玩家，可以用更高频率偷池和施压。面对激进玩家，则要选择更强、更能承受压力的范围继续。

有效激进不是固定风格，而是针对对手漏洞的调整。

\subsection{什么情况下不该盲目激进}

激进有价值，但盲目激进会迅速亏损。以下情况通常不适合强行进攻。

第一，对手不会弃牌。无论你的故事多合理，对手如果拿一对就跟到底，诈唬价值都会大幅下降。

第二，牌面明显有利于对手范围。例如低连接湿润牌面、完成对手主要听牌的转牌、多人底池中的协调牌面。

第三，你的牌没有功能。没有摊牌价值、没有听牌、没有阻断效果、没有后门权益的纯空气牌，在不利牌面上强行下注，通常质量很低。

第四，你没有后续计划。翻牌下注被跟注后，转牌不知道怎么办；转牌下注被跟注后，河牌不知道是否继续。没有计划的进攻容易变成孤立下注。

第五，你是情绪驱动。刚被对手诈唬、刚输大底池、想追回损失、想证明自己，这些都不是进攻理由。

第六，多人底池。多人底池中，对手总范围更容易有人击中强牌，弃牌权益下降，诈唬成功率降低。新手在多人底池中应该显著减少无牌进攻。

激进之前，一定要确认自己不是在用“进攻”掩盖情绪。

\subsection{加注的意义}

加注是比下注更强的进攻行动。它不仅投入更多筹码，也会进一步压缩对手范围。很多新手把加注理解为“我有强牌”，但实际上，加注可以有多种意义。

第一，价值加注。你持有强牌，希望更差牌继续支付。例如你在翻牌中暗三条，面对对手下注，可以加注构建底池。

第二，保护性加注。当你当前领先但牌面湿润，对手可能有许多听牌时，加注可以让对手为继续支付更高价格。不过它仍然不能脱离价值逻辑：你要确认更差牌和听牌会继续。

第三，半诈唬加注。你持有强听牌，通过加注获得弃牌权益。如果对手弃牌，你直接赢；如果对手跟注，你仍然有改善机会。

第四，施压加注。有时你判断对手下注范围较宽，包含大量中等牌力或空气牌，通过加注可以迫使他弃掉一部分范围。

第五，隔离和争取主动权。翻前 3-bet、翻牌加注，都可能让你从被动防守转为主动进攻。

但加注也有风险。它会迅速放大底池，并让对手继续范围变强。如果你用中等牌力随意加注，往往会让更差牌弃牌、更强牌继续。新手使用加注时，必须先明确：我是价值、半诈唬，还是施压？如果说不清，就不要随便加注。

\section{新手常见误区}

\subsection{误区一：只等好牌}

很多新手过于被动，认为只有拿到强牌才能赢。他们错过牌面就放弃，中等牌力只想便宜摊牌，听牌只会跟注等待完成。

这种打法的问题是，赢池方式太单一。你必须依赖发牌帮助，必须摊牌时最好，才能赢底池。对手一旦发现你只在有牌时下注，就可以轻松弃掉弱牌，或者在你过牌时频繁抢池。

德州扑克需要主动争取底池。你不能只等待自己拥有最好牌，也要学会在范围、牌面和对手都合适时，用进攻让对手弃牌。

\subsection{误区二：把激进理解为乱 bluff}

另一类新手走向相反极端。他们听说“要激进”，就开始频繁 bluff，任何牌都想打走对手。结果对手一跟注，他们就没有计划；对手一加注，他们就陷入混乱。

激进不是乱 bluff。激进需要满足条件：范围有优势，牌有功能，对手有可弃牌范围。缺少这些条件，所谓激进只是鲁莽。

真正的激进不是为了显得强，而是为了让对手在错误范围中做困难决策。

\subsection{误区三：该进攻时被动}

有些新手拿到强听牌，却只是被动跟注；处在有利牌面，却因为自己没有成牌而过牌；面对容易弃牌的对手，却不敢开第二枪。

这类错误看似保守，实际上会损失大量 EV。强听牌通过半诈唬可以拥有两个赢法，而被动跟注通常只剩下完成听牌这一条路。有范围优势的牌面，如果你不进攻，就放弃了对对手空气牌和中等牌力施压的机会。

该进攻时被动，会让你变成只能靠摊牌赢的人。

\subsection{误区四：该放弃时硬冲}

也有新手在完全不合适的局面硬冲。牌面明显击中对手，对手是跟注站，自己没有任何权益，却仍然连续下注。这样的打法不是勇敢，而是亏损。

扑克不是每个底池都要争。优秀牌手知道哪些底池值得抢，也知道哪些底池应该放弃。放弃低质量进攻机会，才能保留筹码在高质量机会中施压。

\subsection{误区五：忽略对手类型}

有些玩家使用同一套进攻策略对所有人。对紧弱玩家不敢施压，对跟注站拼命 bluff，对激进玩家用弱牌硬抗。这些都是没有根据对手类型调整的表现。

对手不同，弃牌权益不同。
弃牌权益不同，激进策略就不同。

面对会弃牌的人，进攻可以增加；面对不会弃牌的人，价值下注才是核心。激进不是固定风格，而是针对对手漏洞的策略。

\subsection{误区六：加注没有目的}

新手有时会用加注“看看对手是不是真有牌”，或者因为“不想被欺负”而加注。这和第七章中的试探性下注一样，都是目的不清。

加注必须更谨慎，因为它比普通下注投入更多筹码。加注前要问：

\begin{itemize}
\item 我是价值加注吗？
\item 我是半诈唬加注吗？
\item 我希望对手弃掉哪些牌？
\item 如果对手再加注，我怎么办？
\item 我的牌是否适合打大底池？
\end{itemize}

如果这些问题回答不了，就不要用加注表达情绪。

\section{实战思考框架}

本章的核心框架是“有效激进三条件”：

\begin{center}
\textbf{我的范围有优势。\\
我的牌有合适的功能。\\
对手有足够多可以弃掉的牌。}
\end{center}

\subsection{第一条件：我的范围有优势}

有效进攻首先要看范围优势。你不是因为“想打”而打，而是因为这个牌面、这条行动线或这张转牌更支持你的范围。

例如你翻前 open，大盲跟注，翻牌为 \hand{A♠ 7♣ 2♦ rainbow}。这个牌面通常更有利于翻前加注者，因为你有更多强 A、大对子和高牌组合。你可以更高频率进攻。

再如你翻牌下注被跟注，转牌出现 \hand{K}，而这张 K 明显强化你的翻前强范围。你可以用强价值牌继续，也可以选择合适诈唬牌开第二枪。

相反，如果牌面是 \hand{9♠ 8♠ 7♦ two-tone}、\hand{8♠ 7♣ 6♦}、\hand{T♠ 9♣ 8♦} 这类更容易击中防守方范围的结构，你的范围优势下降，进攻频率应降低。

判断范围优势时，可以问：

\begin{itemize}
\item 这个牌面更容易击中谁的翻前范围？
\item 谁有更多顶对、超对、暗三条、两对、顺子？
\item 这张转牌或河牌强化了谁的故事？
\item 我的前面行动是否能自然代表强牌？
\end{itemize}

没有范围基础的进攻，很容易被对手继续或反击。

\subsection{第二条件：我的牌有合适的功能}

即使范围有优势，也不是所有牌都适合进攻。你还需要判断自己的具体手牌是否有合适功能。

适合进攻的牌通常包括：

\begin{itemize}
\item 强价值牌：希望更差牌跟注；
\item 坚果牌：希望构建大底池；
\item 强听牌：可以半诈唬；
\item 带阻断牌的空气牌：适合选择性诈唬；
\item 有后门权益的牌：被跟注后仍有改善空间；
\item 没有摊牌价值的牌：不下注几乎无法赢，只能选择诈唬或放弃。
\end{itemize}

不适合随意进攻的牌包括：

\begin{itemize}
\item 中等摊牌价值：经常更适合控池；
\item 弱摊牌价值：不值得投入太多筹码；
\item 低质量听牌：被跟注后权益不足；
\item 没有阻断、没有权益、没有故事的空气牌。
\end{itemize}

一手牌是否适合进攻，不只看它强不强，还要看它的功能。强牌负责价值，强听牌负责半诈唬，低摊牌价值空气牌可以选择性诈唬，中等牌力则经常负责防守和摊牌。

\subsection{第三条件：对手有足够多可以弃掉的牌}

进攻能否盈利，最终还要看对手是否会弃牌。

如果对手范围中有大量中等牌力、弱对子、错过听牌和空气牌，并且他有弃牌能力，那么你的进攻会产生弃牌权益。

如果对手范围非常强，或者对手类型是跟注站，那么进攻价值下降。此时你应该减少诈唬，增加价值下注。

判断对手是否适合被施压，可以问：

\begin{itemize}
\item 他翻牌是否跟得太宽？
\item 他转牌和河牌是否会放弃中等牌力？
\item 他是否害怕大底池？
\item 他是否会思考牌面变化？
\item 他是否有过明显 hero call 倾向？
\end{itemize}

对手如果没有弃牌按钮，你就不要用诈唬去寻找他的弃牌按钮。

\subsection{有效激进检查表}

在实战中，你可以使用以下检查表判断是否适合主动进攻：

\begin{enumerate}
\item 我是否拥有范围优势或坚果优势？
\item 我的具体手牌是价值牌、强听牌，还是好诈唬候选？
\item 这个牌面是否支持我的行动线？
\item 对手范围中是否有足够多中等牌力或弱牌？
\item 对手是否有弃牌能力？
\item 如果被跟注，我下一街是否有计划？
\item 如果被加注，我是否知道如何应对？
\end{enumerate}

如果大部分答案为“是”，进攻通常更有依据。
如果大部分答案为“否”，所谓进攻很可能只是鲁莽。

\section{牌例拆解}

\subsection{牌例一：高牌干燥面上的有效进攻}

\subsubsection*{牌局背景}

有效筹码 100BB。你在按钮位拿到 \hand{QJo}，前面玩家弃牌，你 open raise 到 2.5BB。大盲跟注。

翻牌为 \hand{A♠ 7♣ 2♦ rainbow}。大盲过牌，行动轮到你。

\subsubsection*{新手常见想法}

新手可能会想：

\begin{itemize}
\item “我没中牌，应该过牌。”
\item “如果他有 A，我下注就是送钱。”
\item “没有牌不要乱 bluff。”
\end{itemize}

这些想法只看到了自己的具体牌，没有看到范围优势和弃牌权益。

\subsubsection*{理性分析}

第一，范围优势。按钮 open range 中有大量 A、高对和强牌，而大盲防守范围较宽，包含许多没有击中的牌。\hand{A♠ 7♣ 2♦ rainbow} 是典型有利于翻前加注者的高牌干燥面。

第二，手牌功能。你持有 \hand{Q♠ J♣}，没有摊牌价值，属于空气牌。但在这个牌面上，它可以作为小尺度诈唬候选。它不适合摊牌赢，下注至少有机会让更好牌弃牌。

第三，对手弃牌范围。大盲可能持有 K 高、Q 高、小口袋对子、未击中的连张等。这些牌面对小注会有一部分弃牌。

第四，下注尺度。由于牌面干燥，小尺度下注即可发挥作用，不需要大额冒险。

\subsubsection*{推荐行动}

小尺度 C-bet，例如 1/3 底池。

\subsubsection*{本手牌教训}

没中牌不代表不能进攻。只要范围有优势、牌面适合、对手有可弃牌范围，空气牌也可以成为合理诈唬候选。

\subsection{牌例二：低连接面上不该盲目激进}

\subsubsection*{牌局背景}

有效筹码 100BB。你在前位拿到 \hand{AQo}，open raise 到 3BB。大盲跟注。

翻牌为 \hand{9♠ 8♠ 7♦ two-tone}。大盲过牌，行动轮到你。

\subsubsection*{新手常见想法}

一些新手会想：

\begin{itemize}
\item “我是翻前加注者，要持续下注。”
\item “如果我不打，就显得很弱。”
\item “他可能也没中，我可以代表大牌。”
\end{itemize}

这正是把主动权误用为自动进攻。

\subsubsection*{理性分析}

第一，范围优势。这个牌面非常容易击中大盲防守范围。大盲可以有 \hand{T♠ 8♣}、\hand{J♠ T♣}、\hand{9♠ 8♣}、\hand{8♠ 7♣}、\hand{7♠ 6♣}、小对子、两对、暗三条和各种听牌。你的前位 open range 虽然有高对，但在坚果优势上并不明显。

第二，手牌功能。你持有 \hand{AQo}，没有同花听牌，没有顺子听牌，只有两张高牌。它不是强听牌，也不是好的半诈唬候选。

第三，对手弃牌范围。大盲在这个牌面上有大量可继续牌。即使他没有成强牌，也可能有一对、顺子听牌、同花听牌或组合听牌。弃牌权益不足。

\subsubsection*{推荐行动}

倾向过牌。

这不是因为你永远不能下注，而是因为当前三个有效激进条件都不充分：范围优势不足，手牌功能较差，对手可弃牌范围有限。

\subsubsection*{本手牌教训}

主动权不是自动 C-bet 的理由。在低连接湿润面上，没有权益的高牌空气不适合盲目进攻。

\subsection{牌例三：强听牌的半诈唬加注}

\subsubsection*{牌局背景}

有效筹码 100BB。你在按钮位拿到 \hand{A♠ Q♠}，前面玩家弃牌，你 open raise 到 2.5BB。大盲跟注。

翻牌为 \hand{J♠ T♠ 4♦}。大盲主动下注 1/2 底池，行动轮到你。

\subsubsection*{新手常见想法}

新手可能会想：

\begin{itemize}
\item “我还没有成牌，先跟注看转牌。”
\item “加注太激进了，万一他有强牌怎么办？”
\item “等中了同花再打。”
\end{itemize}

跟注当然可以考虑，但这手牌也具备很强的半诈唬加注价值。

\subsubsection*{理性分析}

第一，手牌功能。你拥有坚果同花听牌、两张高牌和顺子相关可能，是非常强的听牌。

第二，弃牌权益。大盲主动下注范围中可能包含 Jx、Tx、弱听牌和部分空气牌。加注可以让其中一部分牌弃牌，尤其是没有能力承受大底池的中等牌力。

第三，被跟注后的权益。即使对手跟注，你仍然有许多转牌改善机会。任何黑桃、K，甚至 A 或 Q 都可能改善你的牌力或带来继续施压机会。

第四，范围表达。你在按钮位 open，可以拥有 \hand{J♠ J♣}、\hand{T♠ T♣}、\hand{J♠ T♣}、强同花听牌和强组合听牌。加注线并非不可信。

\subsubsection*{推荐行动}

可以加注半诈唬，也可以根据对手类型选择跟注。

面对能弃掉中等牌力的对手，加注更有吸引力。面对极度不弃牌但会支付完成听牌的对手，跟注实现权益也合理。

\subsubsection*{本手牌教训}

强听牌不只是用来被动追牌。它们可以通过半诈唬加注同时获得弃牌权益和牌面权益。

\subsection{牌例四：对跟注站不该强行 bluff}

\subsubsection*{牌局背景}

有效筹码 100BB。你在 cutoff 拿到 \hand{Q♠ T♠}，open raise 到 2.5BB。大盲是一名明显跟注站，喜欢用任何对子跟到河牌。

大盲跟注。

翻牌为 \hand{K♠ 7♣ 2♦ rainbow}。大盲过牌，你下注 1/3 底池，大盲跟注。

转牌为 \hand{4♠}。大盲过牌，行动轮到你。你没有成牌，也没有明显听牌。

\subsubsection*{新手常见想法}

一些新手会想：

\begin{itemize}
\item “我翻牌下注了，转牌要继续代表 K。”
\item “他可能只是 7 或小对子，我打大一点能打走。”
\item “如果我放弃，这手牌就输了。”
\end{itemize}

这些想法忽略了最关键的信息：对手类型。

\subsubsection*{理性分析}

第一，范围和牌面。你在 \hand{K♠ 7♣ 2♦} 上翻牌下注有一定合理性，因为高牌干燥面有利于进攻方。但转牌 \hand{4} 是空白牌，并没有显著增加你的故事力量。

第二，手牌功能。你没有成牌，也没有明显听牌或阻断效果。你的牌是低质量空气牌。

第三，对手弃牌能力。大盲是跟注站，翻牌跟注后很可能继续用 Kx、7x、小对子甚至 A 高跟注。你的弃牌权益较低。

第四，继续下注风险。转牌继续下注可能只是扩大亏损。如果他跟注，你河牌大概率仍然没有计划。

\subsubsection*{推荐行动}

过牌放弃。

\subsubsection*{本手牌教训}

对不会弃牌的人，少 bluff。激进必须建立在弃牌权益上。没有弃牌权益的进攻，只是把筹码送给对手。

\subsection{牌例五：价值加注的意义}

\subsubsection*{牌局背景}

有效筹码 100BB。你在大盲拿到 \hand{7♠ 7♣}。按钮位 open raise 到 2.5BB，你跟注。

翻牌为 \hand{K♠ 7♠ 2♦ two-tone}。你过牌，按钮下注 1/3 底池，行动轮到你。

\subsubsection*{新手常见想法}

一些新手会想：

\begin{itemize}
\item “我中了暗三条，跟注让他继续 bluff。”
\item “如果加注，他可能会弃牌。”
\item “慢打可以隐藏牌力。”
\end{itemize}

慢打可以考虑，但不应该成为默认动作。

\subsubsection*{理性分析}

第一，牌力功能。你持有暗三条，是非常强的价值牌，接近坚果。

第二，牌面结构。牌面有两张同花，存在同花听牌。按钮位可能有 Kx、同花听牌、口袋对子和部分空气牌。你有很多目标牌可以继续支付。

第三，加注目的。这是价值加注，同时带有保护意义。你希望 Kx、强听牌和部分不愿弃牌的手牌支付更多筹码。你也希望开始构建大底池，而不是只让对手用小注控制底池。

第四，底池规划。如果只跟注，转牌对手可能过牌，导致你很难在河牌赢到足够多筹码。翻牌加注可以更快建立底池。

\subsubsection*{推荐行动}

加注。

尺度可以根据对手和牌面选择，例如加到对手下注的 3 至 4 倍左右，目标是让 Kx 和听牌继续支付。

\subsubsection*{本手牌教训}

加注不只是 bluff，也可以是价值和构建底池。强牌不能总是慢打，尤其在有听牌和可支付目标时，价值加注非常重要。

\section{本章总结}

德州扑克不是被动等待好牌的游戏。如果你只靠摊牌赢底池，就永远少了一条腿。真正成熟的牌手，既能在摊牌时赢，也能在合适条件下通过进攻让对手弃牌。

本章最重要的概念是主动权和弃牌权益。主动权让你能够率先施压，弃牌权益让你即使当前不是最好牌，也有机会通过下注或加注直接赢下底池。

但激进不是鲁莽。有效激进必须满足三个条件：

\begin{center}
\textbf{我的范围有优势。\\
我的牌有合适的功能。\\
对手有足够多可以弃掉的牌。}
\end{center}

范围优势决定你的故事是否可信。手牌功能决定你的进攻是否有质量。对手可弃牌范围决定你的下注是否真的有弃牌权益。

适合激进的牌通常包括强价值牌、坚果牌、强听牌、带阻断效果或后门权益的诈唬候选。不适合盲目激进的牌包括中等摊牌价值、低质量听牌、没有权益且没有故事的空气牌。

适合进攻的牌面通常是有利于你范围的牌面，例如高牌干燥面、强化你范围的转牌，以及对手范围中有大量中等牌力的局面。不适合盲目进攻的牌面，则包括低连接湿润面、明显完成对手听牌的牌面和多人底池中的复杂牌面。

对手类型决定激进策略的盈利性。面对会弃牌的人，可以合理施压；面对跟注站，应该减少诈唬，增加价值下注；面对激进玩家，要避免用弱牌无计划对抗；面对紧弱玩家，可以增加偷池和施压频率。

加注是更强烈的进攻工具。它可以用于价值、保护、半诈唬和施压，但不能用于试探或情绪表达。加注之前，必须知道自己希望对手做什么，也必须知道面对再加注如何处理。

下一章我们将补充基础计算工具：outs、胜率与底池赔率。只有先理解“这个希望值不值这个价格”，才能在后续章节中更理性地决定何时跟注、何时弃牌。再往后，我们将进入弃牌与跟注，把数学计算与对手范围分析结合起来。

\section{课后训练}

\subsection*{训练一：有效激进三条件检查}

记录最近 10 手你主动下注或加注的牌，逐手回答：

\begin{enumerate}
\item 我的范围是否有优势？
\item 我的牌是否有合适功能？
\item 对手是否有足够多可以弃掉的牌？
\item 如果被跟注，我下一街是否有计划？
\item 如果被加注，我是否知道如何应对？
\end{enumerate}

如果三个核心条件中有两个以上不满足，说明这次进攻可能质量不高。

\subsection*{训练二：找出该进攻却被动的牌}

从最近牌局中找出 5 手你本可以主动进攻、但实际选择过牌或被动跟注的牌。

重点分析：

\begin{itemize}
\item 我是否拥有范围优势？
\item 我的牌是否是强听牌或好诈唬候选？
\item 对手是否容易弃牌？
\item 我当时为什么不敢下注或加注？
\item 如果重新来一次，我会如何设计进攻线？
\end{itemize}

这个训练的目标，是减少过度被动。

\subsection*{训练三：找出不该进攻却硬冲的牌}

从最近牌局中找出 5 手你进攻失败的牌，分析它们是否属于鲁莽进攻。

逐手回答：

\begin{enumerate}
\item 牌面是否真的有利于我的范围？
\item 我的牌是否有牌面权益、阻断效果或后门权益？
\item 对手是否具备弃牌能力？
\item 我是否因为情绪、追损或不想示弱而下注？
\item 我在哪一街应该停止进攻？
\end{enumerate}

这个训练的目标，是区分有效激进和鲁莽进攻。

\subsection*{训练四：半诈唬候选训练}

记录 10 手你持有听牌的牌，判断它们是否适合半诈唬。

重点回答：

\begin{itemize}
\item 这是强听牌还是弱听牌？
\item 被跟注后我有多少改善牌？
\item 我是否有弃牌权益？
\item 我是否有位置？
\item 对手是否会弃牌？
\item 加注、下注和跟注中哪一个更合适？
\end{itemize}

不要把所有听牌都当成半诈唬。强听牌才是高质量进攻候选。

\subsection*{训练五：对手施压对象分类}

在你的常见牌局中，选出 5 名对手，按照以下类型分类：

\begin{itemize}
\item 跟注站；
\item 紧弱玩家；
\item 激进玩家；
\item 松凶玩家；
\item 普通娱乐玩家。
\end{itemize}

对每位玩家回答：

\begin{enumerate}
\item 他是否容易弃牌？
\item 他在哪条街最容易放弃？
\item 他面对大注是否会过度弃牌？
\item 他是否喜欢抓诈唬？
\item 我应该对他增加价值下注，还是增加诈唬？
\end{enumerate}

这个训练的目标，是让激进策略和对手类型匹配。

\subsection*{训练六：加注目的复盘}

记录最近 10 次你加注的牌，逐手回答：

\begin{itemize}
\item 我是价值加注、保护加注、半诈唬加注，还是施压加注？
\item 哪些更差牌会继续？
\item 哪些更好牌会弃牌？
\item 如果对手再加注，我的计划是什么？
\item 这次加注是否让底池大小匹配我的牌力？
\end{itemize}

如果你无法说明加注目的，说明这次加注可能只是情绪化或试探性行动。
